% -----------------------------------------------------------------------------------
\section{Types, parameters, and certification}\label{sec:taxonomy}

A recurring hope in the debate is that the relevant problem types are few: that a decision theory
could be settled by its behaviour on a small catalogue of representative problems. Whether this is
so depends on what counts as the same type, and on which structural axioms a theory is already
known to satisfy.

\subsection{Problems}

\emph{Boolean classifications are automatically finite.} For $k$ local criteria, the profile
$\mathcal S(D)=\{I\subseteq\{1,\dots,k\}:\text{some policy satisfies every criterion in }I\}$ takes at most
$2^{2^k}$ values. A finite classification by joint-satisfiability profiles thus exists trivially; it
need not help to solve any problem, and \cri{4}, being a relation between problems, is not one of
these local predicates.

\emph{Exact recommendation structures are not finite.} Call two problems equivalent if a bijection of
their policy spaces preserves the set of policies acceptable under each criterion.

\begin{proposition}\label{prop:infinitetypes}
Even with two actions and one always-reached information set, the continuous mixed model contains
infinitely many pairwise inequivalent problems.
\end{proposition}

\begin{proof}
For distinct $r_1,\dots,r_n\in(0,1)$, let a reward $Y\in\{0,1\}$, independent of the realised action,
satisfy $\Pr(Y=1\mid p)=1-\prod_i(p-r_i)^2$, a continuous probability since each factor lies in $[0,1]$.
The value is maximised exactly at $p\in\{r_1,\dots,r_n\}$, so exactly $n$ policies satisfy \cri{1}, while
every policy satisfies \cri{2} and \cri{3}. A bijection preserving the \cri{1}-sets cannot identify
problems with different $n$.
\end{proof}

For deterministic policies, Theorem~\ref{thm:realise} already shows that every finite recommendation
structure occurs. The permissive definition of a theory has a further, direct consequence.

\begin{proposition}[No finite benchmark suite]\label{prop:benchmark}
No finite set of problems can establish that an arbitrary decision theory satisfies \cri{1}
everywhere.
\end{proposition}

\begin{proof}
Modify $T^{\mathrm{ex}}$ only on one two-action problem outside the suite, with a strictly better
action, so as to select the worse action. The two theories agree on every benchmark, and only one
satisfies \cri{1} everywhere.
\end{proof}

This holds even if the benchmarks and the omitted problem are of the same intuitive type.
Representative testing requires an invariance requirement on theories. Note also that infinitely
many equivalence classes do not by themselves preclude a finite-dimensional parameterisation: a
single real parameter indexes infinitely many cases.

\subsection{Rules}

The question is sharper for rules on matrices, where invariance axioms are available. Finitely many
parameters must then be read cardinally: a finite-dimensional input space does not make an
arbitrary rule on it a finite-dimensional family.

\begin{proposition}\label{prop:NFO}
There are $2^{\cont}$ rules satisfying \ax{A0}, \ax{N}, \ax{F}, and \ax{O}. Finite benchmark suites cannot
certify \ax{P} or \ax{D} for them.
\end{proposition}

\begin{proof}
For $0<t<1$ let $N_t=\bigl(\begin{smallmatrix}0&-1\\t+1&t\end{smallmatrix}\bigr)$. These lie in distinct
orbits of relabelling combined with positive affine maps, and each has trivial stabiliser: without a
swap, preserving the entries $0$ and $-1$ forces the identity map; with a swap, the diagonal $(t,0)$
must be mapped to $(0,s)$, which requires a negative scale. For each $E\subseteq(0,1)$, choose option $1$
on $N_t$ when $t\in E$ and option $2$ otherwise, transport this ranking equivariantly over each orbit,
and use diagonal ranking elsewhere. This gives $2^{\cont}$ distinct rules, whereas finite tuples of reals
form a set of cardinality $\cont$. In $N_t$, column $1$ strictly dominates column $2$ and the second
diagonal entry is larger; so modifying column sums, or diagonal ranking, on one untested orbit
produces a violation of \ax{P}, or of \ax{D}, that no finite suite (even an adaptive one) detects.
\end{proof}

The construction uses the permission to have arbitrary lookup rules; it does not carry over unchanged
to classes with continuity or measurability requirements.

\begin{proposition}\label{prop:NFB}
There are $2^{\cont}$ rules satisfying \ax{N}, \ax{F}, and \ax{B} (without \ax{O}), and finite suites cannot
certify \ax{D} for them.
\end{proposition}

\begin{proof}
For $E\subseteq\Rl$ and $M_{ii}>M_{jj}$, declare $i\succ_Mj$ if $(M_{ij}-M_{jj})/(M_{ii}-M_{jj})\in E$; never prefer a
lower diagonal strictly to a higher one, and tie all other pairs. The ratio is invariant under positive
affine maps and depends only on the principal block, and strict comparisons descend the diagonal, so
choice sets are nonempty. On $M_t=\bigl(\begin{smallmatrix}1&t\\0&0\end{smallmatrix}\bigr)$ the choice is
$\{1\}$ if $t\in E$ and $\{1,2\}$ otherwise. Starting from $E=\Rl$ (diagonal ranking) and removing one
untested ratio defeats any finite suite.
\end{proof}

\begin{proposition}\label{prop:NFSO}
There are $2^{\cont}$ rules satisfying \ax{N}, \ax{F}, \ax{S}, and \ax{O} that agree on every binary
matrix.
\end{proposition}

\begin{proof}
Let $r_i=(M_{ii}-\min_jM_{ij})/(\max_jM_{ij}-\min_jM_{ij})$, with $r_i=0$ for a constant row, and for
$E\subseteq(0,1)$ rank by $s_i=\one\{r_i\in E\}$. These rules satisfy \ax{N}, \ax{F}, \ax{S}, and \ax{O}, even
under separate positive affine transformations of rows. On binary matrices every $r_i$ is $0$ or $1$,
so every score is zero; on
$\Bigl(\begin{smallmatrix}t&0&1\\1&0&1\\1&1&0\end{smallmatrix}\Bigr)$ the scores are
$(\one\{t\in E\},0,0)$.
\end{proof}

Together with Corollary~\ref{cor:three}, these results give the following picture.

\begin{center}
\small
\renewcommand{\arraystretch}{1.25}
\begin{tabularx}{\textwidth}{@{}>{\raggedright\arraybackslash}p{9em}
                                  >{\raggedright\arraybackslash}X@{}}
  \toprule
  assumptions beyond \ax{A0} & size of the family; certification \\
  \midrule
  \ax{N}, \ax{F}, \ax{O} & $2^{\cont}$; finite suites certify neither \ax{P} nor \ax{D} \\
  \ax{N}, \ax{F}, \ax{B} & $2^{\cont}$; finite suites do not certify \ax{D} \\
  \ax{N}, \ax{F}, \ax{S}, \ax{O} & $2^{\cont}$; all binary behaviour can leave the rule undetermined \\
  \ax{N}, \ax{F}, \ax{C}, \ax{O}, \ax{Q}, \Mon & $2^{\cont}$; finite suites do not certify \ax{D} (Section~\ref{sec:exante}) \\
  \ax{N}, \ax{F}, \ax{B}, \ax{O} & exactly three rules; one unequal-diagonal test distinguishes them \\
  \bottomrule
\end{tabularx}
\end{center}

\enquote{The relevant types are few} is thus true in exactly one of these regimes, and there only
conditionally: when principal binary independence and ordering are known to hold globally, a single
binary test decides among the three possible rules. Finite observation never verifies those
structural axioms themselves.
